\section{Dynamax reuse requires a declared conditional law and information clock}
\label{exec19:dynamax-scope}

A servicing team may infer a borrower's unobserved repayment propensity from
payments, contacts and changing account conditions. A state-space engine is a
natural computational component, but its likelihood must describe the experiment
that generated those observations. Treating an adaptive contact as an ordinary
input can omit information carried by the contact decision itself. Replacing an
online filter by a full-history smoother can introduce information unavailable
when the bank acted. These are distinct scientific errors even when both calls
return finite arrays.

The selected Dynamax revision is commit
\texttt{5098927f2c377c0dd548cda750e6f1436220aac6}, dated 29 September 2026.
Twelve original files at that revision supply the present comparison: the
project description, root notice, dependency metadata, HMM abstraction,
filtering and smoothing implementation, standard transition component,
state-space fitting wrapper, version helper, selected tests and parameter/shape
helpers. This is a historical source inspection. The inherited PyPI version
1.0.2 is a separate release observation; the selected dynamic version helper
has not been executed to establish release equivalence.\citep{exec19_dynamax_source}

\subsection{An input-conditioned model is not a logger-aware causal model}
\label{exec19:dynamax-logger}

The inspected HMM wrapper supplies initial-state probabilities, transition
matrices and emission log likelihoods to its filter. Its marginal-probability
method returns the filter's marginal log likelihood. No separate observed-action
likelihood is automatically inserted. For externally fixed inputs this is a
reasonable computational representation of a declared conditional experiment.
For an action selected using a private state, the joint law generally also needs
the action mechanism. Merely passing the realized action as an input does not
establish that the omitted factor is ignorable.\citep{exec19_dynamax_source}

Consider an authored one-period construction. Let latent repayment propensity
$Z$ be Bernoulli with probability $1/2$. Conditional on $Z$, let action $A$
match $Z$ with probability $9/10$, and outcome $Y$ match $Z$ with probability
$4/5$. The action and outcome noises are conditionally independent. The outcome
has no structural dependence on the action. The observed law therefore factors as
\begin{equation}
 p(z,a,y)=p(z)\,p(a\mid z)\,p(y\mid z).
 \label{exec19:dynamax-law}
\end{equation}
Exact enumeration gives
\begin{equation}
 P(Y=1\mid A=1)=\frac{37}{50},\qquad
 P(Y=1\mid A=0)=\frac{13}{50},\qquad
 E[Y\mid\operatorname{do}(A=1)]-E[Y\mid\operatorname{do}(A=0)]=0.
 \label{exec19:dynamax-contrast}
\end{equation}
The observational contrast is $12/25$, although acting has no effect in the
stipulated structural model. Given $A=1,Y=1$, the posterior probability of
$Z=1$ is $36/37$. Omitting the logger and conditioning only through the outcome
factor gives $4/5$ instead. Both calculations are coherent for their respective
information sets; only the former uses the observed informative action under
this law. They are exact project calculations, not Dynamax execution results.
\citep{exec19_dynamax_reference}

Suppose a unit improvement in repayment probability is assigned USD 200 of
value and an action costs USD 20. These are assumptions. Substituting the
observational contrast gives USD 76, whereas the stipulated causal effect gives
minus USD 20. No model-fitting improvement can make those two targets equal
without additional assumptions or information. A proposed adapter must either
justify the conditional-input experiment, provide the logger-aware law, or
return a clearly described unsupported causal request.\citep{exec19_dynamax_reference}

\subsection{Preserve the observation and transition clocks}
\label{exec19:dynamax-clocks}

The source distinguishes predicted, filtered and smoothed probabilities.
Filtering at time $t$ conditions on the observation prefix; smoothing can use
all observations through the end of the sequence. The distinction belongs in
the product's result type and availability timestamp. An online decision cannot
silently consume the smoother because it is more concentrated or easier to
compute. Retrospective state reconstruction is a legitimate but different task.
\citep{exec19_dynamax_source}

For a separate exact example, let $Z$ persist across two dates, with prior
probability $1/2$, and let two conditionally independent signals each match $Z$
with probability $4/5$. Observing $Y_1=0$ gives
\begin{equation}
 P(Z=1\mid Y_1=0)=\frac15,\qquad
 P(Z=1\mid Y_1=0,Y_2=1)=\frac12.
 \label{exec19:dynamax-filter}
\end{equation}
Suppose a decision between the two signals earns USD 100 when $Z=1$ and costs
USD 30. Under that explicitly assumed response rule, the legitimate first-date
expected net is minus USD 10. Using the later signal would report USD 20 for
the same first-date decision. This is an information-clock violation, not an
observed advantage of smoothing or evidence about real servicing returns.
\citep{exec19_dynamax_reference}

Input indexing is equally consequential. The generic HMM transition helper
constructs its time-varying transition matrices from \texttt{inputs[1:]},
consistent with the inspected model abstraction's destination-time input.
The low-level filter description instead writes its transition using the
source-time input. A wrapper must declare the mapping rather than infer it from
an argument name. The standard transition class accepts an input argument but
returns its fixed transition matrix without using that argument. Passing a
contact sequence therefore does not ensure that the selected transition model
actually responds to contacts. A deterministic indexed-sequence reference
should compare sampling, model wrappers and low-level calls before admission.
\citep{exec19_dynamax_source}

\subsection{Numerical completion and parameter learning need separate checks}
\label{exec19:dynamax-numerics}

The inspected normalizer preserves exact zeros, floors very small positive
entries and substitutes one for a zero denominator. An all-zero input consequently
remains all zero rather than becoming a valid probability distribution. The
emission update subtracts the maximum log likelihood without a dedicated
all-minus-infinity branch. The smoother masks near-zero predicted probabilities
using a closeness test before another normalization. These are local source
observations motivating native conformance cases; they are not measured defect
rates or a claim about every reachable backend execution. Invalid support,
nonfinite arithmetic, unsupported laws and valid negative scientific findings
must not collapse into the same output status.\citep{exec19_dynamax_source}

The fitting wrapper exposes a specified iteration count and an objective trace.
For nonpositive iteration counts it returns the supplied parameters with an
empty trace. The reviewed path does not make a general convergence or
identification guarantee. Selected parameter helpers also distinguish fully
frozen groups from partially frozen groups that certain joint EM updates
reject; the inspected tests compare that boundary with SGD support. A caller
must retain initialization, trainable flags, shapes, optimization settings and
failures. A state posterior conditional on fitted parameters is not automatically
a confidence statement about parameter estimates or a selected policy.
\citep{exec19_dynamax_source}

Selected upstream tests compare filtering and smoothing with explicit finite
latent-path sums, exercise nonstationary transitions and inspect fitting
behavior. They provide useful candidate conformance specifications, but none
was executed here. A future native suite should include impossible emissions,
zero support, one-step and empty sequences, input shifts, frozen parameters,
multiple starts and device synchronization. Resource and timing claims must
include compilation, conversion and failed requests.\citep{exec19_dynamax_source}

\subsection{Prefer a narrow original-engine comparison; defer executable admission}
\label{exec19:dynamax-choice}

The selected root notice is MIT. The project metadata requires Python at least
3.10 and eight direct runtime dependencies: JAX, JAXlib, TFP-nightly,
fastprogress, Optax, jaxtyping, typing-extensions and NumPy. Those declarations
are unversioned; build and optional groups are separate. A complete distribution,
transitive notices and resolved backend/precision environment remain necessary
before executable admission. Root license text alone does not qualify the whole
installed stack.\citep{exec19_dynamax_source}

The proposal should retain the original engine as a conditional-state-space
comparator behind a narrow adapter. A small exact enumerator remains the
transparent baseline for finite cases; it does not offer a scaling solution.
A complete port would duplicate numerical, differentiation and maintenance
obligations without a demonstrated platform need. A causal policy optimizer
built solely from a conditional likelihood is outside the initial profile.
Scientific ownership must cover the target, logger and clocks; engineering
ownership must cover the environment, conversions and failure contract; release
ownership must cover the actual distribution and notices.

This is a bounded reuse recommendation with executable admission deferred.
Public sources and generated examples suffice to specify the first comparison.
Native conformance, parameter uncertainty and measured full-workflow cost remain
prototype work; authentic bank assignment and value remain later institutional
questions. Neither the package name nor an API returning a posterior closes
those distinct obligations.
